\documentclass[12pt,a4paper]{article}

% --- GÓI TIẾNG VIỆT VÀ ĐỊNH DẠNG TRANG ---
\usepackage[utf8]{inputenc}
\usepackage[vietnamese]{babel}
\usepackage[T5]{fontenc}
\usepackage[top=2cm, bottom=2cm, left=3cm, right=2cm]{geometry} % Chuẩn lề đồ án
\usepackage{mathptmx} % Font Times New Roman tương thích pdfLaTeX
\usepackage{setspace}
\onehalfspacing % Giãn dòng 1.5 lines

% --- CÁC GÓI HỖ TRỢ BẢNG, HÌNH VÀ VIỀN ---
\usepackage{tikz} % Vẽ khung viền trang bìa
\usetikzlibrary{calc}
\usepackage{booktabs}
\usepackage{tabularx}
\usepackage{array}
\usepackage{multirow}
\usepackage[hidelinks]{hyperref} % Liên kết mục lục tự động, không hiện khung đỏ

\begin{document}

% =========================================================================
% 1. TRANG BÌA (COVER PAGE)
% =========================================================================
\begin{titlepage}
    \begin{tikzpicture}[remember picture, overlay]
        % Khung viền đôi chuẩn đồ án
        \draw[line width=3pt, color=black] 
            ($(current page.north west) + (2.5cm,-1.5cm)$) 
            rectangle 
            ($(current page.south east) + (-1.5cm,1.5cm)$);
        \draw[line width=1pt, color=black] 
            ($(current page.north west) + (2.6cm,-1.6cm)$) 
            rectangle 
            ($(current page.south east) + (-1.4cm,1.4cm)$);
    \end{tikzpicture}

    \begin{center}
        \vspace*{0.2cm}
        {\large \textbf{TRƯỜNG ĐẠI HỌC SÀI GÒN}}\\[0.2cm]
       
        \textbf{BỘ MÔN HỆ THỐNG THÔNG TIN}\\[0.5cm]
        
        \rule{0.5\textwidth}{1pt}\\[1.5cm]

        {\large \textbf{BÁO CÁO ĐỒ ÁN}}\\[0.3cm]
        {\large \textbf{PHÂN TÍCH THIẾT KẾ HỆ THỐNG THÔNG TIN}}\\[1.2cm]

        % TÊN ĐỀ TÀI
        {\Large \textbf{ĐỀ TÀI:}}\\[0.4cm]
        {\LARGE \textbf{XÂY DỰNG HỆ THỐNG THÔNG TIN}}\\[0.3cm]
        {\LARGE \textbf{QUẢN LÝ KHO SÁCH}}\\[0.4cm]
        {\large \textbf{(TRIỂN KHAI PHÂN HỆ B: NHẬP XUẤT KHO)}}\\[2.0cm]
    \end{center}

    \vfill

    \begin{flushleft}
        \hspace*{1.0cm}
        \begin{tabular}{ll}
            \textbf{Giảng viên hướng dẫn:} & Lê Nhị Lãm Thúy \\[0.4cm]
            \textbf{Nhóm sinh viên thực hiện:}  \\[0.2cm]
            1. Lê Phước Đại & MSSV:  \\
            2. Phạm Nguyễn Tuấn Kiệt & MSSV: [MSSV Kiệt] \\
            3. Huỳnh Hữu Phúc & MSSV: [MSSV Hữu Phúc] \\
            4. Nguyễn Minh Tâm & MSSV: [MSSV Tâm] \\
            5. [Họ và tên Hoàng Phúc] & MSSV: [MSSV Hoàng Phúc] \\
        \end{tabular}
    \end{flushleft}

    \vspace{1.0cm}
    \begin{center}
        \textbf{TP. HỒ CHÍ MINH, THÁNG [..] NĂM 2026}
    \end{center}
\end{titlepage}

% =========================================================================
% ĐÁNH SỐ TRANG LA MÃ CHO PHẦN ĐẦU (FRONT MATTER)
% =========================================================================
\newpage
\pagenumbering{roman}
\setcounter{page}{1}

% =========================================================================
% 2. BẢNG KẾ HOẠCH LÀM VIỆC VÀ TỰ ĐÁNH GIÁ TỶ LỆ % ĐÓNG GÓP
% =========================================================================
\section*{BẢNG KẾ HOẠCH LÀM VIỆC VÀ TỰ ĐÁNH GIÁ TỶ LỆ ĐÓNG GÓP}
\addcontentsline{toc}{section}{Bảng kế hoạch làm việc và tự đánh giá tỷ lệ đóng góp}

\noindent\textbf{1. Nguyên tắc đánh giá:}
\begin{itemize}
    \item Thành viên có tỷ lệ đóng góp cao nhất được tính là 100\% (hưởng trọn số điểm đồ án).
    \item Điểm số của các thành viên còn lại được tính tỷ lệ thuận theo \% đóng góp đã được nhóm thống nhất.
\end{itemize}

\vspace{0.3cm}

\noindent\textbf{2. Bảng tổng hợp phân công và kết quả thực hiện:}

\vspace{0.2cm}
\noindent
\begin{table}[h!]
\centering
\small
\begin{tabularx}{\textwidth}{|c|>{\raggedright\arraybackslash}p{2.8cm}|c|>{\raggedright\arraybackslash}X|>{\centering\arraybackslash}p{1.8cm}|c|}
\hline
\textbf{STT} & \textbf{Họ và tên} & \textbf{MSSV} & \textbf{Nhiệm vụ phân công chính} & \textbf{Kết quả hoàn thành} & \textbf{Tỷ lệ (\%)} \\
\hline
1 & [Họ và tên Đại] & [MSSV] & Khảo sát nghiệp vụ; Phân tích yêu cầu vào/ra; Thiết kế CSDL (lược đồ quan hệ, mô tả 7 bảng); Lập bảng kế hoạch. & Hoàn thành 100\% & ...\% \\
\hline
2 & [Họ và tên Kiệt] & [MSSV] & Vẽ hệ thống sơ đồ (BFD, Ngữ cảnh, DFD mức 1, 2); Thiết kế giao diện tổng thể; Hướng dẫn cài đặt/sử dụng; Tổng kết; Ghép báo cáo. & Hoàn thành 100\% & ...\% \\
\hline
3 & [Họ và tên Hữu Phúc] & [MSSV] & Lập trình phân hệ B3 (Nhập/xuất kho, kiểm tra tồn kho, bộ lọc giá); Soạn kịch bản và biên bản kiểm thử B3. & Hoàn thành 100\% & ...\% \\
\hline
4 & [Họ và tên Tâm] & [MSSV] & Lập trình phân hệ B4 (Thêm/sửa/xóa Sách, NCC, lịch sử nhập kho, in PDF); Viết thuyết minh B4 và biên bản kiểm thử. & Hoàn thành 100\% & ...\% \\
\hline
5 & [Họ và tên Hoàng Phúc] & [MSSV] & Lập trình phân hệ Admin (Quản trị user, phân quyền); Tích hợp CSDL mẫu; Chuẩn bị tệp SQL, readme.txt và bản demo. & Hoàn thành 100\% & ...\% \\
\hline
\end{tabularx}
\end{table}

\vspace{0.8cm}
\noindent
\begin{minipage}[t]{0.5\textwidth}
    \centering
    \textbf{Ý KIẾN CỦA GIẢNG VIÊN}\\
    \textit{(Ký và ghi rõ họ tên)}
\end{minipage}
\begin{minipage}[t]{0.5\textwidth}
    \centering
    \textbf{ĐẠI DIỆN NHÓM SINH VIÊN}\\
    \textit{(Ký và ghi rõ họ tên)}\\[2cm]
    [Họ tên người đại diện nhóm]
\end{minipage}

% =========================================================================
% 3. MỤC LỤC TỰ ĐỘNG, DANH MỤC HÌNH ẢNH VÀ BẢNG BIỂU
% =========================================================================
\newpage
\tableofcontents % Tự động sinh mục lục

\newpage
\listoffigures % Tự động sinh danh mục hình ảnh
\addcontentsline{toc}{section}{\listfigurename}

\newpage
\listoftables % Tự động sinh danh mục bảng biểu
\addcontentsline{toc}{section}{\listtablename}

% =========================================================================
% BẮT ĐẦU VÀO PHẦN NỘI DUNG CHÍNH (ĐÁNH SỐ TRANG 1, 2, 3...)
% =========================================================================
\newpage
\pagenumbering{arabic}
\setcounter{page}{1}

% --- Ví dụ các chương tiếp theo sẽ được Kiệt ghép vào đây: ---
% \include{chuong1_gioithieu}
% \include{chuong2_phantich}
% \include{chuong3_thietke}
% \include{chuong4_caidat_sudung}
% \include{chuong5_tongket}

\end{document}
